\documentclass[12pt]{article}
\usepackage{amsmath,amssymb,booktabs}
\usepackage[utf8]{inputenc}
\usepackage[T2A]{fontenc}
\usepackage[russian,english]{babel}
\title{Транс-реактивная сеть как синтетическая агентная модель информационной фрагментации}
\author{TRN Synthetic Modeling Package}
\date{}
\begin{document}
\maketitle

\section{Аннотация}
Рассматривается синтетическая агентная модель TRN как параметризованной информационной среды. Модель предназначена для анализа консенсуса, поляризации, фрагментации и антиконсенсуса. Реальные персональные данные и реальные платформенные API не используются.

\section{Модель агентов}
Каждый агент задаётся вектором
\[
x_i(t)=(b_i(t),e_i(t),q_i(t),r_i(t),m_i(t)).
\]

\section{Динамика}
\[
b_i(t+1)=clip_{[-1,1]}[b_i(t)+\Delta t(S_i(t)+I_i(t)+\eta_i(t))].
\]

\section{TRN-поле}
\[
I_i(t)=\lambda m_i(1-r_i)(1-q_i)A_i(t)(P_i(t)-b_i(t)).
\]

\section{Антиконсенсус}
Антиконсенсус фиксируется при
\[
C(t)<0.45,\quad Pol(t)>0.44,\quad Ext(t)>0.35.
\]

\section{Индекс риска}
\[
R_{TRN}=\frac{\lambda\bar m(1-\bar r)(1-\bar q)\chi}{\bar h+\epsilon}.
\]

\end{document}
